\chapter{Implementazione}
\label{cap:implementazione}
L'implementazione deve rappresentare una simulazione deterministica. Per questo motivo non sono stati usati thread (oggetti Java \texttt{Thread}), la gestione del tempo avviene in modo discreto ed è centralizzata nello scheduler.

%%%%%%%%%%%%%%%%%%%%%%%%%%%%%%%%%%%%%%%%%%%%%%%%%%%%%%%%%%%%
\section{Scheduler}
Ogni possibile implementazione di un algoritmo di scheduling è un'estensione della classe \texttt{Scheduler}. Questo permette di definire i comportamenti comuni a tutti gli scheduler e di astrarre future implementazioni.

Per l'implementazione è stato necessario mantenere una lista di job pronti \texttt{readyTasks} e una di job bloccati \texttt{blockedTask}. Inoltre è necessario mantenere un riferimento all'ultimo job che è andato in esecuzione.

\myskip

L'implementazione della logica di scheduling è la stessa sia per RM che per EDF: il metodo \texttt{schedule} con i metodi helper necessari sono stati, per questo motivo, inseriti all'interno della classe base. Ogni implementazione concreta deve solo definire alcuni aspetti usati poi nella logica. Per modellare questo principio è stato usato il Template Method pattern: il metodo \texttt{schedule} è dichiarato pubblico e \texttt{final} in modo che non possa essere modificato dalle implementazioni, mentre gli hooks chiamati al suo interno sono dichiarati \texttt{abstract} e \texttt{protected}.

Per gestire i tempi rilevanti, cioè quelli in cui lo scheduler deve prendere il controllo del sistema per rivedere la propria politica di scheduling, è delegata alla struttura \texttt{readyTasks}. Gli eventi rilevanti sono l'unione ordinata dei multipli di ciascun periodo fino alla durata della simulazione.

Per capire i passaggi che esegue la simulazione consideriamo e analizziamo il sequence diagram di Figura~\ref{fig:sequenceDiagram}:
\begin{itemize}
    \item \texttt{assignPriority} \\
        Assegna le priorità nominali ai task. È dichiarato astratto in \texttt{Scheduler} e deve essere implementato dalle classi concrete.
    \item \texttt{MyClock.reset} \\
        Resetta il clock di sistema prima di iniziare la simulazione.
    \item \texttt{initStructure} \\
        Inizializza le strutture dati: la lista di job pronti e la mappa dei job attivi per task. Rilascia i primi job al tempo 0.
    \item \texttt{checkFeasibility} \\
        Controlla la schedulabilità del taskset.
    \item \texttt{executeFor} \\
        Metodo dello scheduler che manda in esecuzione il job a priorità maggiore per un intervallo di tempo. Gestisce l'avanzamento del clock e il logging degli eventi di esecuzione, preemption e completamento.
    \item \texttt{Job.execute} \\
        Riduce il tempo di esecuzione rimanente del job.
    \item \texttt{releasePeriodTasks} \\
        Si occupa di rilasciare i nuovi job allo scoccare del periodo di ogni task e controlla che il job precedente abbia terminato, altrimenti solleva una \texttt{DeadlineMissException}.
\end{itemize}

\begin{figure}[htbp]
    \centering
    \includegraphics[width=.65\textwidth]{immagini/sequence diagram.pdf}
    \caption{Sequence Diagram scheduling}
    \label{fig:sequenceDiagram}
\end{figure}

\subsection{Rate Monotonic}
Vediamo come sono implementati i metodi astratti della classe \texttt{Scheduler} in RM.
\begin{itemize}
    \item \texttt{checkFeasibility} \\
        Viene controllato l'hyperbolic bound.
    \item \texttt{assignPriority} \\
        Assegna la priorità in modo inverso rispetto alla durata del periodo.
    \item \texttt{addReadyTask} \\
        Aggiunge il job alla lista ordinata secondo la priorità dinamica.
\end{itemize}

\subsection{Earliest Deadline First}
Vediamo come sono implementati i metodi astratti della classe \texttt{Scheduler} in EDF.
\begin{itemize}
    \item \texttt{checkFeasibility} \\
        Viene controllato che il fattore di utilizzo del taskset sia minore o uguale a 1.
    \item \texttt{assignPriority} \\
        Assegna la priorità in modo inverso rispetto alla durata della deadline.
    \item \texttt{addReadyTask} \\
        Aggiunge il job alla lista e la riordina secondo la deadline assoluta del job.
\end{itemize}

%%%%%%%%%%%%%%%%%%%%%%%%%%%%%%%%%%%%%%%%%%%%%%%%%%%%%%%%%%%%
\section{Fault injection}
Il sistema permette di modellare fault legati al tempo di esecuzione. Essendo il tempo di esecuzione campionato da una distribuzione fornita al task, è possibile modellare comportamenti stocastici che superano il caso medio o il WCET previsto.

%%%%%%%%%%%%%%%%%%%%%%%%%%%%%%%%%%%%%%%%%%%%%%%%%%%%%%%%%%%%
\section{Utility}
\subsection{Logging}
Il logging è centralizzato nello scheduler e scrive su un file una sequenza di coppie $<tempo, evento>$. Il file di destinazione è \texttt{trace.log}.

\subsection{Clock}
\label{subsec:clock}
Il tempo è gestito globalmente tramite il Singleton \texttt{MyClock}. Lo scheduler è l'unica entità responsabile di far avanzare il clock tramite i metodi \texttt{advanceTo} e \texttt{advanceBy}. Il sistema gestisce il tempo in nanosecondi per garantire precisione, esponendolo in millisecondi all'utente.

\subsection{Eccezioni}
\begin{itemize}
    \item \texttt{DeadlineMissedException} \\
        Viene sollevata dallo scheduler quando un job supera la propria deadline o quando un nuovo job viene rilasciato prima che il precedente sia terminato.
\end{itemize}
